\documentclass[11pt,a4paper]{article}
\usepackage[T1]{fontenc}
\usepackage[utf8]{inputenc}
\usepackage{lmodern}
\usepackage[czech]{babel}
\usepackage[margin=2cm]{geometry}
\usepackage{amsmath,amssymb,mathtools}
\usepackage{xcolor,cancel,enumitem}
\usepackage{tikz,pgfplots}
\pgfplotsset{compat=1.18}
\usetikzlibrary{arrows.meta,calc,patterns,babel}
\pagestyle{empty}
\setlength{\parindent}{0pt}
\setlength{\parskip}{4pt}
\newcommand{\vysl}[1]{\boxed{\color{red!75!black}#1}}
\newcommand{\pozn}[1]{{\color{red!75!black}#1}}
\newcommand{\dx}{\,\mathrm{d}x}
\newcommand{\dt}{\,\mathrm{d}t}
\newcommand{\subst}[1]{\left|\begin{array}{l}#1\end{array}\right|}
\DeclareMathOperator{\arctg}{arctg}
\DeclareMathOperator{\arccotg}{arccotg}
\DeclareMathOperator{\tg}{tg}
\DeclareMathOperator{\cotg}{cotg}
\usepackage{tcolorbox}
\newtcolorbox{pomocny}{colframe=red!75!black, colback=white, coltitle=white, fonttitle=\bfseries, title={Pomocný výpočet:}, sharp corners}
\newcommand{\dy}{\,\mathrm{d}y}
% znaménkové schéma: \znam{bod/popis/typ,...}{znaménka v intervalech}
% typ: o = vyloučený bod (kroužek), x = obyčejný bod
\newcommand{\znam}[2]{%
\begin{tikzpicture}[baseline=-3pt]
\def\n{0}
\foreach \p/\l/\t [count=\i] in {#1} {\xdef\n{\i}}
\pgfmathsetmacro{\L}{1.6*(\n+1)}
\draw (0,0) -- (\L,0);
\foreach \p/\l/\t [count=\i] in {#1} {
  \pgfmathsetmacro{\xx}{1.6*\i}
  \ifnum\pdfstrcmp{\t}{o}=0 \draw[fill=white] (\xx,0) circle (2.5pt); \draw (\xx-0.09,0) -- (\xx+0.09,0) (\xx,-0.09) -- (\xx,0.09);
  \else \draw (\xx,0.1) -- (\xx,-0.1); \fi
  \node[below=2pt] at (\xx,0) {\small $\l$};
}
\foreach \s [count=\i] in {#2} {
  \pgfmathsetmacro{\xx}{1.6*\i-0.8}
  \node[above] at (\xx,0) {$\s$};
}
\end{tikzpicture}}
\newcommand{\aug}[1]{\left(\begin{array}{rrr|rrr}#1\end{array}\right)}
\newcommand{\krok}[1]{\text{\small\color{orange!80!black}#1}}

\begin{document}
\section*{Určete $D(f)$ a graf:}

\textbf{\textcircled{1}}\quad $y=\ln(x+3)+1$

$\left.\begin{array}{r} x+3>0\\ x>-3\end{array}\right\}\ D(f)=(-3,\infty)$

\begin{minipage}{0.55\textwidth}
\begin{tikzpicture}
\begin{axis}[axis lines=middle, width=8cm, height=6.5cm, clip=true,
  xmin=-3.6, xmax=4.5, ymin=-2.6, ymax=3.6, xtick={-3,-2,1}, ytick={1}, xlabel=$x$, ylabel=$y$]
\addplot[thick, blue!70!black, domain=0.08:4.4, samples=80] {ln(x)} node[pos=0.9, above] {$\ln x$};
\addplot[thick, green!45!black, domain=-2.92:2, samples=80] {ln(x+3)} node[pos=0.75, below right] {$\ln(x+3)$};
\addplot[thick, red!80!black, domain=-2.97:0.5, samples=80] {ln(x+3)+1};
\node[red!80!black, rotate=40] at (axis cs:-1.6,2.6) {$\ln(x+3)+1$};
\draw[dashed] (axis cs:-3,-2.6) -- (axis cs:-3,3.6);
\draw[red!80!black, -{Stealth}] (axis cs:1,0) -- (axis cs:-2,0);
\draw[red!80!black, -{Stealth}] (axis cs:-2,0) -- (axis cs:-2,1);
\draw[dashed, gray] (axis cs:-2,1) -- (axis cs:0,1);
\end{axis}
\end{tikzpicture}
\end{minipage}%
\begin{minipage}{0.42\textwidth}
\begin{enumerate}[label=\textcircled{\arabic*}]
\item $y=\ln x$
\item {\color{green!45!black}$y=\ln(x+3)$} (posun o 3 vlevo)
\item {\color{red!80!black}$y=\ln(x+3)+1$} (posun o 1 nahoru)
\end{enumerate}
\end{minipage}

\hrule
\medskip
\textbf{\textcircled{2}}\quad $y=\left|\sqrt{x+4}-2\right|$

$\left.\begin{array}{r} x+4\ge0\\ x\ge-4\end{array}\right\}\ D(f)=\langle-4,\infty)$

\begin{minipage}{0.55\textwidth}
\begin{tikzpicture}
\begin{axis}[axis lines=middle, width=8.5cm, height=6.5cm, clip=false,
  xmin=-4.6, xmax=7, ymin=-2.6, ymax=3.6, xtick={-4,1,4}, ytick={1,2}, xlabel=$x$, ylabel=$y$]
\addplot[thick, blue!70!black, domain=0:7, samples=80] {sqrt(x)} node[above] {$\sqrt x$};
\addplot[thick, green!45!black, domain=-4:2.5, samples=80] {sqrt(x+4)} node[right] {$\sqrt{x+4}$};
\addplot[thick, brown!60!black, domain=-4:6.5, samples=80] {sqrt(x+4)-2} node[below right] {$\sqrt{x+4}-2$};
\addplot[thick, red!80!black, domain=-4:6.5, samples=80] {abs(sqrt(x+4)-2)};
\node[red!80!black, rotate=25] at (axis cs:-3.2,2.6) {$\left|\sqrt{x+4}-2\right|$};
\draw[dashed, gray] (axis cs:-4,2) -- (axis cs:0,2);
\draw[dashed, gray] (axis cs:4,0) -- (axis cs:4,2);
\end{axis}
\end{tikzpicture}
\end{minipage}%
\begin{minipage}{0.42\textwidth}
\begin{enumerate}[label=\textcircled{\arabic*}]
\item $y=\sqrt x$
\item {\color{green!45!black}$y=\sqrt{x+4}$} (posun o 4 vlevo)
\item {\color{brown!60!black}$y=\sqrt{x+4}-2$} (posun o 2 dolů)
\item {\color{red!80!black}$y=\left|\sqrt{x+4}-2\right|$} záporné hodnoty se \uv{překlopí} kolem osy $x$ do kladných
\end{enumerate}
\end{minipage}

\newpage
\textbf{\textcircled{3}}\quad $y=2-\dfrac{x^3}{2}$, \qquad $D(f)=\mathbb{R}$

\begin{minipage}{0.5\textwidth}
\begin{tikzpicture}
\begin{axis}[axis lines=middle, width=7cm, height=9cm, clip=true,
  xmin=-2.6, xmax=2.6, ymin=-5, ymax=5.5, xtick={-2,-1,1,2}, ytick={-4,-2,-1,1,2,4}, ticklabel style={font=\scriptsize}]
\addplot[thin, blue!70!black, domain=-2:2, samples=60] {x^3};
\addplot[thick, green!45!black, domain=-2.3:2.3, samples=60] {x^3/2};
\addplot[thin, brown!60!black, domain=-2.3:2.3, samples=60] {-x^3/2};
\addplot[thick, red!80!black, domain=-2.2:2.2, samples=60] {2-x^3/2};
\node[blue!70!black] at (axis cs:1.95,5) {$x^3$};
\node[green!45!black] at (axis cs:2.3,3.6) {$\frac{x^3}{2}$};
\node[brown!60!black] at (axis cs:2.15,-4.2) {$-\frac{x^3}{2}$};
\node[red!80!black] at (axis cs:-1.5,4.8) {$2-\frac{x^3}{2}$};
\end{axis}
\end{tikzpicture}
\end{minipage}%
\begin{minipage}{0.48\textwidth}
\begin{enumerate}[label=\textcircled{\arabic*}]
\item {\color{blue!70!black}$y=x^3$}
\item {\color{green!45!black}$y=\frac{x^3}{2}$} (hodnoty na ose $y$ se zmenší na polovinu)
\item {\color{brown!60!black}$y=-\frac{x^3}{2}$} (hodnoty se změní na opačné -- překlopení kolem $x$)
\item {\color{red!80!black}$y=2-\frac{x^3}{2}$} (posun o 2 nahoru)
\end{enumerate}
\end{minipage}

\hrule
\section*{Určete $D(f)$:}

\textbf{\textcircled{1}}\quad $y=\sqrt{\dfrac{x-2}{x+3}}$

\begin{center}
\begin{tikzpicture}[scale=1.2]
\node at (0.3,1.7) {$\dfrac{x-2}{x+3}\ge0$};
\node (p0) at (5.3,1.7) {pro $x=0 \to \dfrac{0-2}{0+3}=-\dfrac23<0$};
\draw[-{Stealth}, semithick] ($(p0.south)+(0.3,0.05)$) to[out=-100, in=30] (3.3,0.65);
\draw (-1,0) -- (7,0);
\draw[very thick, red!80!black] (-1,0) -- (1.5,0);
\draw[very thick, red!80!black] (4.5,0) -- (7,0);
\draw[fill=white] (1.5,0) circle (3pt); \draw (1.38,0) -- (1.62,0) (1.5,-0.12) -- (1.5,0.12);
\fill (4.5,0) circle (2.5pt);
\node at (0.3,0.45) {\textcircled{$+$}}; \node at (3,0.45) {\textcircled{$-$}}; \node at (5.8,0.45) {\textcircled{$+$}};
\node[below, align=center] at (1.5,-0.1) {$-3$\\ \small$(x+3=0$\\ \small$x=-3)$\\ \small jmenovatel};
\node[below, align=center] at (4.5,-0.1) {$2$\\ \small$(x-2=0$\\ \small$x=2)$\\ \small čitatel};
\end{tikzpicture}
\end{center}
\[
\vysl{D(f)=(-\infty,-3)\cup\langle2,\infty)}
\]

\newpage
\textbf{\textcircled{2}}\quad $y=\ln\dfrac{x^2+5x-14}{x^2+x-6}$
\[
\frac{x^2+5x-14}{x^2+x-6}>0, \qquad\text{upravme:}\quad \frac{(x-2)(x+7)}{(x+3)(x-2)}>0
\]
\begin{center}
\begin{tikzpicture}
\draw (0,0) -- (6.4,0);
\draw (1.6,0.1) -- (1.6,-0.1);
\foreach \xx in {3.2,4.8} {\draw[fill=white] (\xx,0) circle (2.5pt); \draw (\xx-0.09,0) -- (\xx+0.09,0) (\xx,-0.09) -- (\xx,0.09);}
\node[below=2pt] at (1.6,0) {\small $-7$};
\node[below=2pt] at (3.2,0) {\small $-3$};
\node[below=2pt] at (4.8,0) {\small $2$};
\node[above] at (0.8,0) {$+$};
\node[above] at (2.4,0) {$-$};
\node[above] (pm) at (4.0,0) {\textcircled{$+$}};
\node[above] at (5.6,0) {$+$};
\node[anchor=north west] (p0) at (-0.6,-0.85) {pro $x=0$: \quad $\dfrac{(0-2)(0+7)}{(0+3)(0-2)}=\dfrac{-14}{-6}>0$};
\draw[-{Stealth}, semithick] ($(p0.north east)+(-0.75,-0.1)$) to[out=75, in=0] (pm.east);
\end{tikzpicture}
\end{center}

\underline{Pozor:} $x\ne2$ \dots\ aby ve jmenovateli nebyla 0, ale lze krátit $(x-2)$ $\Rightarrow$ neovlivní znaménko
\[
\vysl{D(f)=(-\infty,-7)\cup(-3,2)\cup(2,\infty)}
\]

\hrule
\medskip
\textbf{\textcircled{3}}\quad $y=\ln(\ln x)$

\begin{minipage}{0.55\textwidth}
\begin{align*}
x &> 0\\
\ln x &> 0\\
x &> 1
\end{align*}
\[
\vysl{D(f)=(1,\infty)}
\]
\end{minipage}%
\begin{minipage}{0.4\textwidth}
\begin{tikzpicture}
\begin{axis}[axis lines=middle, width=5cm, height=4cm, clip=false,
  xmin=-0.5, xmax=4, ymin=-1.5, ymax=1.8, xtick={1}, ytick=\empty, xlabel=$x$, ylabel=$y$]
\addplot[thick, domain=0.25:3.8, samples=60] {ln(x)} node[above] {$\ln x$};
\addplot[line width=2.2pt, red!80!black, opacity=0.75, domain=1:3.8, samples=60] {ln(x)};
\addplot[only marks, mark=*, mark size=1.8pt, red!80!black] coordinates {(1,0)};
\end{axis}
\end{tikzpicture}\\
z grafu je vidět, že pro $x>1$ je $\ln x>0$
\end{minipage}

\end{document}
